\documentclass[10pt]{article}

\usepackage[utf8]{inputenc}

\usepackage[T1]{fontenc}

\usepackage{amsmath}

\usepackage{amsfonts}

\usepackage{amssymb}

\usepackage[version=4]{mhchem}

\usepackage{stmaryrd}

\usepackage{mathrsfs}



\begin{document}

\begin{enumerate}

  \item $\Phi_{0}$ - тождественное отображение;



  \item $\Phi_{t}$ - нормальное вполне положительное отображение такое, что $\Phi_{t}(I)=I$ для любого $t>0$;



  \item $\Phi_{t} \Phi_{s}=\Phi_{t+s}$ для $t, s \geqslant 0$;



  \item $\Phi_{t}(X) \rightarrow X$ при $t \rightarrow 0$ для любого $X \in d$ в ультраслабой операторной топологии.



\end{enumerate}



Если .d - коммутативная алгебра функций на нек-ром «фазовом пространстве» $\mathscr{E}$, то это соответствует полугруппе переходных операторов в теории марковских случайных процессов; если же . операторов в гильбертовом пространстве $H$, то говорят о квантовой динамической полугруппе.



Важной проблемой математич. теории Д. П. является характеризация инфинитезимального оператора $\mathscr{\varphi}$, к-рый определяется формулой



\[

\mathscr{L}(X)=\lim _{t \rightarrow 0+} t^{-1}\left[\Phi_{t}(X)-X\right]

\]



где имеется в виду предел в ультраслабой операторной топологии. Если область определения $D(\mathscr{L})$ есть *-подалгебра .d, то $\mathscr{\rho}$ является вполне диссипативным отображением таким, что $\varphi(I)=0$ (см. Диссипативное отображение). Инфинитезимальные операторы полностью описаны лишь для непрерывных по норме Д. п. В этом случае $\varphi$ является ограниченным отображением в .d вида



\[

\mathscr{L}(X)=i[H, X]+\Phi(X)-\frac{1}{2}\{\Phi(I) X+X \Phi(I)\}, \quad X \in \mathscr{A},

\]



где $\Phi$ - вполне положительное отображение в .ø, $\mathrm{H}$ - самосопряженный оператор из ./ (см. [2]). Для квантовой Д. П. это сводится к формуле Линдблада:



\[

\mathscr{L}(X)=i[H, X]+\sum_{j}\left\{L_{j}^{*} X L_{j}-\frac{1}{2}\left(L_{j}^{*} L_{j} X+X L_{j}^{*} L_{j}\right)\right\},

\]



где $H=H^{*}, L_{j}, \Sigma L_{j}^{*} L_{j} \in \mathfrak{B}(H)(\mathrm{cм}$. [7], [5]).



В физич. приложениях квантовые Д. п. возникают как решения марковского управляющего уравне н и я



\[

\frac{d}{d t} \Phi_{t}(X)=\mathscr{P}\left(\Phi_{t}(X)\right)

\]



к-рое получается при рассмотрении взаимодействия квантовой открытой системы с резервуаром в марковском приближении (см. [1], [3]). Формулы (1), (2) дают наиболее обший вид марковского управляющего уравнения, согласующийся с вероятностной интерпретацией квантовой механики. В конкретных моделях операторы $H, L_{j}$ в соотношениях (1), (2) определяются гамильтонианом взаимодействия и состоянием резервуара. Эргодич. свойства Д. П. описывают в марковском приближении процесс релаксации открытой системы к равновесному состоянию (см. [4]).



В общем плане естественна постановка обратной проблемы р а с ш и р е н и я Д. П. до группы автоморфизмов, то есть представления марковской динамики квантовой системы в виде проекции обратимой динамики системы, взаимодействующей с окружением (см. [6]). Возможность такого расширения, означающая согласованность понятия Д. П. с принципами статистич. механики, обуславливается, главным образом, свойством полной положительности.



Jum.: [1] Alick i R., Le ndi K., "Lect. Notes in Phys.», 1987, v. 286; [2] Christensen E., Evans D. E., "J. Lond. Math. Soc.", 1979, v. 20, p. 358-68; [3] Davies E. B., Quantum theory of open systems, L.- N.Y., 1976; [4] Frige ri o A., «Commun. Math. Phys.», 1978, v. 63, p. 269-76; [5] G or in i V. [a. o.], «Rep. Math. Phys.», 1978, v. 13, p. 149-73; [6] Ku m merer B., «J. Funct. Anal.», 1985, v. 63, p. 139-77; [7] Lindblad G., "Comm. Math. Phys.», 1979, v. 65, $\begin{array}{ll}\text { p. 281-94. } & \text { А. С. Холево. }\end{array}$



ДИНАМИЧЕСКАЯ СИММЕТРИЯ Квантовой системЫ, алгебра, генерируюшая спектр, группане и н в а р и а т ности, - симметрия полного пространства векторов состояния системы, образующих одно неприводимое представление некоторой группы или алгебры Ли, операторы которой объединяют в одно семейство все состоя- ния системы и включают в себя операторы переходов между различными состояниями.



Вырождение уровней энергии квантовой системы, находящейся в стационарном состоянии, связано с наличием у нее нек-рой симметрии (группы инвариантности), то есть с наличием набора операторов, коммутирующих с гамильтонианом системы, к-рые обычно образуют конечномерную алгебру Ли. Помимо вырождений, связанных с явной симметрией гамильтониана (напр., относительно вращений в трехмерном пространстве), существует скрытая симметрия, объясняюшая так наз. случайное вырождение уровней энергии системы. Примером такой симметрии, объясняющей вырождение уровней с одинаковым главным квантовым числом и различными орбитальными моментами в атоме водорода, является симметрия $O(4)$ в импульсном пространстве (фоко в ская с иммет р ия; предложена В. А. Фоком в 1935). Аналогично «случайное» вырождение уровней трехмерного изотропного гармонич. осциллятора связано с наличием у него симметрии относительно унитарной группы $U(3)$. Операторы алгебры соответствуюших групп переводят одно выбранное состояние, принадлежащее заданному уровню энергии, во все остальные состояния, принадлежащие тому же уровню энергии; при этом ортогональные состояния, принадлежашие данному уровню, образуют базис неприводимого представления группы симметрии (группы инвариантности).



В отличие от группы инвариантности, действие операторов динамич. группы (группы неинвариантности, или динамич. алгебры Ли) на одно выбранное стационарное состояние квантовой системы порождает все остальные стационарные состояния системы, связывая таким образом все стационарные состояния системы, в том числе принадлежащие различным уровням, в одно семейство - м ульти пле т. При этом группа симметрии (группа инвариантности) системы является подгруппой группы Д.с. Так, для атомов водорода группой Д. с. является конформная $O(4,2)$ динамич. группа, одно неприводимое вырожденное представление к-рой содержит все его связанные состояния, а для трехмерного квантового гармонич. осциллятора - группа $U(3,1)$. Среди генераторов группы Д. с. обязательно есть не коммутирующие с гамильтонианом, действие к-рых переводит волновые функции состояний с одним уровнем энергии квантовой системы в волновые функции состояний с другими энергиями (то есть соответствует квантовым переходам между уровнями системы).



Нахождение динамич. группы симметрии физич. задачи, с одной стороны, эквивалентно решению Шредингера уравнения (или Дирака уравнения, Клейна - Гордона уравнения) для данной системы, с другой стороны - позволяет использовать хорошо развитый математич. аппарат теории представлений групп Ли и получать соотношения типа рекуррентных соотношений для матричных элементов операторов физич. величин, что важно при расчетах физич. эффектов по теории возмущений (напр., при расчете эффекта Штарка для атома водорода).



Группа Д.с. Квантовой системы определяется неоднозначно. Так, для атома водорода наряду с конформной группой $O(4,2)$ Д. с. может являться также группа де Ситтера $O(4,1)$, а для трехмерного осциллятора - неоднородная симплектич. группа $\operatorname{ISp}(6, R)$ [для $N$-мерного осциллятора - ISp $(2 N, R)]$. Выбор той или иной группы Д. с. квантовой системы определяется удобством при расчетах.



В физике элементарных частиц интерес к Д. с. связан с попытками установить симметрию лагранжиана взаимодействия по известному из опыта спектру масс частиц.



Лит.: [1] Малк и н И. А., Ман ько В. И., Динамические симметрии и когерентные состояния квантовых систем, М., 1979.



В. И. Манько.



ДИНАМИЧЕСКАЯ СИММЕТРИЯ т вердого те ла свойство тела, заключающееся в том, что эллипсоид инерции является эллипсоидом врашения (при этом симметричность самого тела не обязательна).



ДИНАМИЧЕСКАЯ СИСТЕМА - (в первоначальном значении термина) механическая система с конечным числом





\end{document}